\documentclass[twoside,11pt]{article}
\usepackage[hmargin=1in,vmargin=1in]{geometry}
\usepackage{../style/jeffe,graphicx}
\usepackage[utf8x]{inputenc}
\usepackage{microtype}
\usepackage{cite}
\usepackage{enumerate}
\usepackage{asymptote}


\usepackage[charter]{mathdesign}
\usepackage[scaled=.96,osf]{XCharter}
\usepackage{sourcesanspro,inconsolata,eucal}
\SetMathAlphabet{\mathsf}{bold}{\encodingdefault}{\sfdefault}{b}{\updefault}
\SetMathAlphabet{\mathtt}{bold}{\encodingdefault}{\ttdefault}{b}{\updefault}
\SetMathAlphabet{\mathsf}{normal}{\encodingdefault}{\sfdefault}{\mddefault}{\updefault}
\SetMathAlphabet{\mathtt}{normal}{\encodingdefault}{\ttdefault}{\mddefault}{\updefault}
\usepackage{textcomp}

\newtheorem{lemma}{Lemma}[section]
\newtheorem{theorem}[lemma]{Theorem}
\newtheorem{corollary}[lemma]{Corollary}
\newtheorem{conjecture}[lemma]{Conjecture}
\numberwithin{figure}{section}

\def\Sym#1{\texttt{#1}}						% symbol (generator/terminal/dart)
\def\Inverse#1{\={#1}}						% inverse

\def\rev{\operatorname{\mathit{rev}}}		% reversal
\def\Head{\operatorname{\mathit{head}}}		% head
\def\Tail{\operatorname{\mathit{tail}}}		% tail
\def\Left{\operatorname{\mathit{left}}}		% left shore
\def\Right{\operatorname{\mathit{right}}}	% right shore
\def\bdry{\partial\!}

\makeatletter
\newcommand{\smash@bar}[4]{%
  \smash{\rlap{\raisebox{-#3\fontdimen5#10}{$\m@th#2\mkern#4mu\mathchar'26$}}}%
}
\newcommand{\lambdabar}{{\mathchoice
  {\smash@bar\textfont\displaystyle{0.25}{2.75}\lambda}
  {\smash@bar\textfont\textstyle{0.25}{2.75}\lambda}
  {\smash@bar\scriptfont\scriptstyle{0.25}{2.75}\lambda}
  {\smash@bar\scriptscriptfont\scriptscriptstyle{0.25}{2.75}\lambda}
}}
\makeatother

\def\Walk{\omega}
\def\PerturbedWalk{\widetilde\omega}
\def\Radius{\rho}
\def\Relevant#1{\overline{#1}}
\def\Lift#1{\widehat{#1}}
\def\Tiling{\Lift{Q}}
\def\Grammar{\mathcal{G}}

\usepackage{arydshln}
\dashlinedash 0.75pt
\dashlinegap 1.5pt

\clubpenalty 8000
\widowpenalty 8000

\let\Newpage\relax

% =========================================================

\begin{document}

\pagestyle{myheadings}
\markboth{Topologically Trivial Closed Walks in Directed Surface Graphs}
    {Jeff Erickson and Yipu Wang}

\begin{titlepage}

\title{Static Trees%
\thanks{Research on this paper was partially supported by NSF grant CCF-1408763.  An extended abstract of this paper will be presented at the 33rd International Symposium on Computational Geometry \cite{ccwx}.}}

\author{\href{http://jeffe.cs.illinois.edu}{Jorge Velez}\\[1ex]
University of Illinois at Urbana-Champaign}

\maketitle
It is my responsibilty as a human being that has acquired knowledge to try to understand \emph{reality} and seek \emph{truth} from the existance of what I percieve and yet not perceive based
on mathematical \emph{proof}. Anything outside of this horizon is irrelevant to my existance it is \emph{meaningless} in a reality embodied by \emph{Nothingness}. 
\end{titlepage}

\tableofcontents

\section{P vs. NP Statement}
An array is a linear data structure that collects elements of the same data type and stores them in contiguous memory locations." \\ 
To begin with we can acknowledge that we know far more today about \emph{computional complexity} than we knew before starting with \emph{Kurt Gödel} and \emph{John Von Neumann}. 
But for what we know about the subject we still are far from truly understanding what encompasses computional complexity. \\

So we being with a general \textbf{Statement:} Given some map with labeled countries is there a configuration such that the map is 3-colorable?  The questions we should generally ask ourselves is the following 
is there an efficient algorithm to compute and tells us if a given map is 3-colorable? And also we must define what me mean by effiecent? So now I claim that effiecnt means that with a resonable amount of steps we 
can get an answer for the 3-colorable problem, where steps is how many computaions are preformed by a given \emph{machine}, where the \emph{machine} here is defined by a \emph{Turing machine}. So now if this \emph{machine} can
compute the 3-clorable problem in \emph{polynomial} time then \textbf{P = NP} and the universe will colopase to those who understand its significance. \\

So more formally we define \textbf{Efficent} as 
\[
    \text{Efficent: } \text{poly}(n)-\text{time} (n = \vert \text{ input } \vert) \:: (P)
 \]

 We also have in all cases where the task has a brute force solution that leads to \emph{exponential}($n$) time algorithm. In the 3-colorable problem we could try all possibel colorings for a given map $M$ that will explode exppnentailly over each try.

\end{document}